% Acronyms
\newacronym[description={\glslink{apig}{Application Program Interface}}]
    {api}{API}{Application Program Interface}

\newacronym[description={\glslink{umlg}{Unified Modeling Language}}]
    {uml}{UML}{Unified Modeling Language}

\newacronym[description={\glslink{spag}{Single Page Application}}]
    {spa}{SPA}{Single Page Application}

\newacronym[description={\glslink{aig}{Artificial Intelligence}}]
    {ai}{AI}{Artificial Intelligence}
    
\newacronym[description={\glslink{erpg}{Enterprise Resource Planning}}]
    {erp}{ERP}{Enterprise Resource Planning}

\newacronym[description={\glslink{crmg}{Customer Relationship Management}}]
    {crm}{CRM}{Customer Relationship Management}
    
\newacronym[description={\glslink{rfidg}{Radio-Frequency Identification}}]
    {rfid}{RFID}{Radio-Frequency Identification}

\newacronym[description={\glslink{llmg}{Large Language Model}}]
    {llm}{LLM}{Large Language Model}

\newacronym[description={\glslink{pocg}{Proof of Concept}}]
    {poc}{PoC}{Proof of Concept}

\newacronym[description={OpenTelemetry}]
    {otel}{OTel}{OpenTelemetry}

\newacronym[description={\glslink{sdkg}{Software Development Kit}}]
    {sdk}{SDK}{Software Development Kit}

\newacronym[description={\glslink{crudg}{Create, Read, Update, Delete}}]
    {crud}{CRUD}{Create, Read, Update, Delete}

\newacronym[description={\glslink{mcpg}{Model Context Protocol}}]
    {mcp}{MCP}{Model Context Protocol}

% Glossary entries
\newglossaryentry{apig} {
    name=\glslink{api}{API},
    text=Application Program Interface,
    sort=api,
    description={in informatica con il termine \emph{Application Programming Interface API} (ing. interfaccia di programmazione di un'applicazione) si indica ogni insieme di procedure disponibili al programmatore, di solito raggruppate a formare un set di strumenti specifici per l'espletamento di un determinato compito all'interno di un certo programma. La finalità è ottenere un'astrazione, di solito tra l'hardware e il programmatore o tra software a basso e quello ad alto livello semplificando così il lavoro di programmazione}
}

\newglossaryentry{umlg} {
    name=\glslink{uml}{UML},
    text=UML,
    sort=uml,
    description={in ingegneria del software \emph{UML, Unified Modeling Language} (ing. linguaggio di modellazione unificato) è un linguaggio di modellazione e specifica basato sul paradigma object-oriented. L'\emph{UML} svolge un'importantissima funzione di ``lingua franca'' nella comunità della progettazione e programmazione a oggetti. Gran parte della letteratura di settore usa tale linguaggio per descrivere soluzioni analitiche e progettuali in modo sintetico e comprensibile a un vasto pubblico}
}

\newglossaryentry{spag} {
    name=\glslink{spa}{SPA},
    text=SPA,
    sort=spa,
    description={con \emph{SPA, Single Page Application} si indica un'applicazione web costituita da una singola pagina HTML che viene aggiornata dinamicamente senza ricaricare completamente il contenuto. Le SPA migliorano l’esperienza utente rendendo la navigazione più fluida e veloce, grazie al caricamento asincrono dei dati e alla gestione lato client della logica applicativa}
}

\newglossaryentry{aig} {
    name=\glslink{ai}{AI},
    text=AI,
    sort=ai,
    description={con \emph{AI, Artificial Intelligence} si indica un insieme di tecniche e metodi informatici che consentono a un sistema di simulare capacità tipiche dell’intelligenza umana, quali apprendimento, ragionamento, percezione e risoluzione di problemi}
}

\newglossaryentry{erpg} {
    name=\glslink{erp}{ERP},
    text=ERP,
    sort=erp,
    description={con \emph{ERP, Enterprise Resource Planning} si indica un sistema informativo aziendale integrato utilizzato per gestire e automatizzare i principali processi operativi di un’organizzazione, quali contabilità, gestione delle risorse umane, produzione, logistica e vendite}
}

\newglossaryentry{crmg} {
    name=\glslink{crm}{CRM},
    text=CRM,
    sort=crm,
    description={con \emph{CRM, Customer Relationship Management} si indica un sistema o insieme di strumenti software utilizzati per gestire le relazioni e le interazioni con clienti attuali e potenziali. Un CRM consente di organizzare e analizzare informazioni relative ai clienti, supportare le attività di vendita, marketing e assistenza, e migliorare la qualità del servizio e la fidelizzazione}
}

\newglossaryentry{rfidg} {
    name=\glslink{rfid}{RFID},
    text=RFID,
    sort=rfid,
    description={con \emph{RFID, Radio-Frequency Identification} si indica una tecnologia che utilizza onde radio per identificare e tracciare oggetti, persone o animali attraverso dispositivi chiamati tag RFID. Questi tag contengono informazioni che possono essere lette a distanza da appositi lettori, senza necessità di contatto diretto o visibilità, rendendo la tecnologia particolarmente utile in ambiti come logistica, tracciamento e gestione degli inventari}
}

\newglossaryentry{llmg} {
    name=\glslink{llm}{LLM},
    text=LLM,
    sort=llm,
    description={con \emph{LLM, Large Language Model} si indica un modello di intelligenza artificiale basato su reti neurali di grandi dimensioni, addestrato su ampie quantità di dati testuali, in grado di comprendere e generare linguaggio naturale}
}

\newglossaryentry{scrumg} {
    name=SCRUM,
    text=SCRUM,
    sort=scrum,
    description={metodologia agile per la gestione e lo sviluppo di progetti, basata su un approccio iterativo e incrementale. Il lavoro viene suddiviso in cicli temporali chiamati Sprint, al termine dei quali viene prodotto un incremento funzionante del sistema. Scrum definisce ruoli (come Scrum Master e Product Owner), eventi (Sprint, Daily Scrum, Sprint Review e Retrospective) e artefatti (Product Backlog e Sprint Backlog) per favorire collaborazione, adattabilità e miglioramento continuo del processo di sviluppo}
}

\newglossaryentry{pocg} {
    name=\glslink{poc}{PoC},
    text=Proof of Concept,
    sort=poc,
    description={con \emph{PoC, Proof of Concept} si indica un prototipo o implementazione preliminare sviluppata per verificare la fattibilità di una soluzione tecnica o di un’idea progettuale. Un PoC ha lo scopo di validare determinati aspetti, come prestazioni, integrazione o corretto funzionamento, prima di procedere con uno sviluppo completo}
    }

\newglossaryentry{sdkg} {
    name=\glslink{sdk}{SDK},
    text=Software Development Kit,
    sort=sdk,
    description={con \emph{SDK, Software Development Kit} si indica insieme di strumenti, librerie, documentazione e API messi a disposizione per facilitare lo sviluppo di applicazioni su una determinata piattaforma o tecnologia. Un SDK fornisce componenti riutilizzabili e funzionalità predefinite che permettono agli sviluppatori di integrare rapidamente servizi specifici, riducendo tempi e complessità di sviluppo.}
    }

\newglossaryentry{crudg} {
    name=\glslink{crud}{CRUD},
    text={Create, Read, Update, Delete},
    sort=crud,
    description={con \emph{CRUD, Create, Read, Update, Delete} si indica l'insieme delle quattro operazioni fondamentali per la gestione dei dati in un sistema informatico. Rappresenta le azioni di creazione, lettura, aggiornamento ed eliminazione di entità all’interno di un database o di un’applicazione. Queste operazioni costituiscono la base per la maggior parte delle funzionalità di persistenza e manipolazione dei dati.}}
    
\newglossaryentry{mcpg} {
    name=\glslink{mcp}{MCP},
    text={Model Context Protocol},
    sort=mcp,
    description={\emph{MCP, Model Context Protocol} è un protocollo che consente a un agente basato su modelli di intelligenza artificiale di interagire con sistemi e servizi esterni in modo strutturato, permettendo l'accesso a dati dinamici e contestuali. MCP abilita l’integrazione tra modelli linguistici e strumenti operativi, migliorando la capacità decisionale dell’agente e riducendo la dipendenza da informazioni statiche}
}
        